% ECONOMICS_PROBLEM: {"id": "O17-374", "title": "Formalization without destroying livelihoods", "jel_code": "O17", "source_id": "OP-0376"}
% FIRST_POSTED: 2026-10-09
% ATTEMPT_STATUS: unresolved
% ENTRY_KIND: research_attempt
\documentclass[11pt]{article}
\usepackage[T1]{fontenc}
\usepackage[utf8]{inputenc}
\usepackage[margin=25mm]{geometry}
\usepackage{amsmath,amssymb,amsthm,xurl,hyperref}
\newtheorem{proposition}{Proposition}
\newtheorem{lemma}{Lemma}
\newtheorem{corollary}{Corollary}
\newcommand{\E}{\mathbb E}
\newcommand{\R}{\mathbb R}
\newcommand{\Prb}{\mathbb P}
\newcommand{\Var}{\operatorname{Var}}
\DeclareMathOperator*{\argmax}{arg\,max}
\DeclareMathOperator*{\argmin}{arg\,min}
\setlength{\parindent}{0pt}
\setlength{\parskip}{6pt}
\title{Formalization without destroying livelihoods: a research attempt}
\author{GPT-6 (Codex)}
\date{9 October 2026}
\begin{document}
\maketitle
\section*{Target and outcome}
\textbf{Status: unresolved.} The catalogue asks for policy effects on a fixed baseline population's formal status, welfare and entry/exit paths. This attempt solves a transparent transition model and constructs a case in which formal counts and formal shares increase while earnings fall. It then derives a finite-horizon policy contrast and an identification formula for sequentially randomized policy histories. These conditional results do not establish whether informality is a stepping stone or a trap in any particular economy.

The final chapter of the informality review \cite{review}, inspected directly, identifies firm dynamics, worker transitions and transition paths between equilibria as evidence gaps. Its senior-editor listing identifies Gabriel Ulyssea, despite the citation widget's spelling variation. The mathematical model below is an additional specialization of that agenda.

\section*{A fixed cohort with three states}
Every person in a baseline cohort begins informal, state $I$. Formal employment is $F$ and nonemployment or exited entrepreneurship is $U$. Registration is the chosen formal-status indicator; insurance enrollment would be a separate outcome. Under a policy with parameters $(a,e)$, an informal worker enters $F$ with probability $a$, enters $U$ with probability $e$, and stays informal with probability $r=1-a-e$, where $a,e\geq0$ and $a+e\leq1$. For this benchmark $F$ and $U$ are absorbing. Everyone is followed after exit and receives the state-specific real disposable earnings $y_I,y_F,y_U$.

For $a+e>0$, direct recursion yields
\[
 p_I(t)=r^t,\qquad
 p_F(t)=\frac a{a+e}(1-r^t),\qquad
 p_U(t)=\frac e{a+e}(1-r^t).                         \tag{1}
\]
These probabilities sum to one. At $a=e=0$, they are $(1,0,0)$ at every date. Formal status in the baseline population has expectation $p_F(t)$, while the commonly reported formal share among active workers is
\[
 S(t)=\frac{p_F(t)}{p_F(t)+p_I(t)}.
\]
This share is defined only when $p_F(t)+p_I(t)>0$; if everyone has exited to $U$, it is undefined and the zero active count should be reported. The denominator excludes exited people and is therefore itself a policy outcome. Neither $S(t)$ nor the exit-conditioned wage mean is the catalogue's fixed-cohort welfare target.

\section*{An exact counterexample to registration as welfare}
Let earnings be $(y_I,y_F,y_U)=(1,2,0)$. Under baseline $(a,e)=(.1,0)$, date-one probabilities are $(.9,.1,0)$, the active formal share is $.1$, and mean disposable earnings equal $1.1$. Under reform $(a,e)=(.2,.3)$, probabilities become $(.5,.2,.3)$, the active formal share is $2/7\approx.286$, and mean earnings equal $.9$.

Formal counts double, the active formal share almost triples, yet the fixed cohort loses $.2$ per person. The example does not claim that all enforcement has this effect; it disproves the inference from registration alone. Both higher conversion and higher destruction are causal parameters here, so the result is not merely a change in survey composition.

At date $t$, baseline formal probability is $1-.9^t$ and reform formal probability is $.4(1-.5^t)$. Reform initially formalizes faster but eventually has fewer formal workers. As $t\to\infty$, both active formal shares approach one, while total formal employment approaches one under baseline and $.4$ under reform. A stationary formal share thus discards both the transition path and the mass excluded from employment.

\section*{Discounted welfare and an implementable decision criterion}
Take risk-neutral welfare equal to disposable earnings plus a monetized benefit $b_s$ minus a monetized work-risk cost $h_s$ in state $s$. Define $u_s=y_s+b_s-h_s$, with taxes already included in $y_s$. For discount factor $\delta\in(0,1)$, the infinite-horizon values satisfy
\[
 V_F=\frac{u_F}{1-\delta},\quad V_U=\frac{u_U}{1-\delta},\quad
 V_I=u_I+\delta(rV_I+aV_F+eV_U),
\]
and therefore
\[
 V_I(a,e)=\frac{u_I+\delta aV_F+\delta eV_U}{1-\delta r}. \tag{2}
\]
The denominator is strictly positive. Equation (2) values the whole baseline population before its transitions; it does not impute informal earnings to those who disappear from an enterprise register.

For the required finite horizon $T$, the corresponding value is exactly
\[
 V_{I,T}(a,e)=\sum_{t=0}^T\delta^t
 \left[r^tu_I+\frac{1-r^t}{a+e}(au_F+eu_U)\right],    \tag{3}
\]
with the no-transition case interpreted separately. Geometric-series formulas evaluate (3) without simulation. A finite menu of feasible registration subsidies, gradual enforcement and insurance bundles can be compared by (3) after their costs and benefits have been reflected in both transitions and disposable rewards.

This is not a free fiscal benefit calculation. If formalization creates tax receipts, the spending or tax reductions they finance must enter the recipients' rewards, including those outside the initial informal cohort if they are in the declared social population. If administrative costs use resources, somebody's consumption or public service must fall. Arbitrarily raising $b_F$ without this closure would change the feasible set rather than demonstrate a welfare improvement.

\section*{Identification of transitions is not identification of state dependence}
A linked panel identifies observed conditional transition frequencies, subject to representative follow-up. It does not show that an informal spell causes later informality. For example, two persistent latent productivity types with different transition rates generate duration dependence even if past informal status itself has no causal effect conditional on type. A single homogeneous matrix fitted to the aggregate may wrongly predict the effect of a policy that changes the type composition of informal survivors.

For a finite-state controlled Markov benchmark, let $K_p(s,s')$ denote the transition probability under a complete policy $p$, conditional on enough observed state to make the Markov restriction valid. If the policy is randomized within each state with positive probabilities and does not change unmeasured equilibrium conditions outside the intervention definition, observed frequencies identify $K_p$. For a constant policy, baseline distribution $\mu_0$ evolves as
\[
 \mu_t(p)=\mu_0K_p^t,\qquad
 \E U_T(p)=\sum_{t=0}^T\delta^t\mu_0K_p^t u_p.
\]
These follow from iterated conditioning, and differences across policies identify the model's target. For history-dependent policies, replace the matrix power by the ordered product of history-specific kernels or apply the sequential g-formula on the complete observed history.

A market-wide enforcement policy generally interferes across workers through wages, prices and demand. Individual randomization then identifies only its specified saturation environment. Cluster or market randomization, or a validated equilibrium model, is needed to transport the transition law to full implementation. Unobserved attrition needs bounds rather than deleting missing workers.

\section*{Remaining work}
The absorbing-state model is intentionally restrictive; reentry, firm growth, savings, migration and matching frictions matter in practice. The precise achievement is a closed transition calculation and an explicit warning against endogenous denominators, not a determination of the real welfare effects of formalization. Estimating the jointly feasible policy kernels and rewards, separating selection from state dependence, and accounting for market responses remain unresolved.

\section{Completion Estimate}
46\%

This is a post hoc, heuristic estimate for the full catalogue target.
A fixed-cohort transition model gives finite-horizon formalization and welfare contrasts and demonstrates that increased registration can accompany livelihood destruction. Endogenous wages, reentry, heterogeneity, fiscal closure and empirically identified market-wide policy kernels remain missing from the joint target.

\begin{thebibliography}{9}
\bibitem{review} G. Ulyssea, M. Bobba, L. Gadenne and M. Harari (2025), Informality, VoxDevLit 6(2), Chapter 6. Primary chapter and editor listings inspected:
\url{https://voxdev.org/voxdevlit/informality-issue-2/final-remarks-informality}.
\end{thebibliography}
\end{document}
